\documentclass [prd,reprint,eqsecnum,nofootinbib]  {revtex4-1}
\usepackage{amsfonts,amsmath,amscd,mathrsfs,hyperref}
\usepackage{xcolor} % Dodano pakiet do obsługi kolorów
\def\SO{\mathsf{SO}}
\def\so{\mathsf{so}}
\def\SU{\mathsf{SU}}
\def\cP{{\cal P}}
\def\cD{{\cal D}}
\def\cC{{\cal C}}
\def\cL{{\cal L}}
\def\w{\omega}
\def\v{\wedge}

\begin{document}
\title{%
Local Maxwell symmetry and gravity}
\author{\textcolor{blue}{A. C. Vance}}
\email{a.vance@arkham-physics.edu}\affiliation{\textcolor{blue}{Department of Theoretical Physics, Arkham State University, Miskatonic Valley, 02-100 Arkham, USA}}
\author{\textcolor{blue}{E. R. Hawthorne}}
\email{e.hawthorne@arkham-physics.edu}\affiliation{\textcolor{blue}{Department of Theoretical Physics, Arkham State University, Miskatonic Valley, 02-100 Arkham, USA}}
\date{\today}

\begin{abstract}\noindent
\textcolor{blue}{In this manuscript, we investigate} some physical aspects of a theory obtained
by gauging the AdS-Maxwell symmetry. \textcolor{blue}{Such a} theory has the form of
Einstein gravity coupled to the $\sf{SO(3,1)}$ Yang-Mills field. We
notice that there is another tetrad field, which can be associated
with a linear combination of Lorentz and Maxwell connections. Taking
this tetrad as a fundamental variable makes it possible to cast the
theory into the form of $f-g$ gravity in \textcolor{blue}{a} first order formulation.
Finally, \textcolor{blue}{we extend the previous analysis by presenting a detailed cosmological model and solving the remaining analytical dynamics of the auxiliary fields.}
\end{abstract}
\maketitle

\section{Introduction}

The Maxwell symmetry is an extension of Poincar\'e symmetry, being a
symmetry of fields on constant electromagnetic background
\cite{Bacry:1970ye}, \cite{Schrader:1972zd}. This symmetry is of
interest from a purely algebraic point of view because it circumvents
a well-known theorem that does not allow for central extension of
this algebra (see e.g., \cite{Galindo:1967}, \cite{Soroka:2011tc},
\cite{Gomis:2009dm}).

However, until recently, the Maxwell symmetry did not attract much
interest. This is \textcolor{green}{somehow} surprising, because physical systems
living on constant electromagnetic background, or closely related to such a
situation, are frequently encountered and important in
physics. This suggests that the Maxwell symmetry may play a quite
important role in physics. For example, since canonical
non-commutativity (for review see \cite{Douglas:2001ba}) is
ultimately related to the behavior of quantum systems on constant
electromagnetic background (or generalization thereof), one can
expect that this symmetry may play an important role in
non-commutative field theories (as already noticed in
\cite{Doplicher:1994tu}). By the same token, it does not seem
excluded that it might be relevant in the context of the quantum Hall
effect.

Leaving these speculations to future works, in this paper we
investigate the related problem, namely the construction and
properties of theories with local Maxwell invariance. In the next
section, following \cite{Durka:2011nf}, we recall the formulation of
such a theory, deriving also its field equations. The following
short section is devoted to some comments on possible relations
between the gauged Maxwell theories and bi-metric theories of
gravity. \textcolor{blue}{Section IV is dedicated to cosmological solutions, while Section V provides a novel analytical and numerical treatment of the auxiliary field dynamics.}

\section{Gauged Maxwell theory}

We begin this section \textcolor{blue}{by} presenting some definitions and then recalling
some of the results presented in \cite{Durka:2011nf}.

The starting point of our construction is a generalization of
Maxwell symmetry, in which the Poincar\'e subalgebra of the original
Maxwell algebra \cite{Bacry:1970ye}, \cite{Schrader:1972zd} is
replaced with the anti-de Sitter counterpart. In other words, we
consider the algebra of symmetries of (complex scalar, Dirac) fields
on anti-de Sitter space and \textcolor{blue}{a} constant electromagnetic
background\footnote{A detailed discussion will be presented in a
forthcoming paper.}. The algebra of such symmetries, which turns out
to be isomorphic to $\sf{so}(3,1)\oplus\sf{so}(3,2)$, was discussed
previously in \cite{Soroka:2006aj}, \cite{Bonanos:2010fw},
\cite{Lukierski:2010dy} and has the following form
\cite{Durka:2011nf}:
\begin{align}
&[\mathcal{P}_{a},\mathcal{P}_{b}]=i(\mathcal{M}_{ab}-\mathcal{Z}_{ab})\,
,\nonumber\\
&[\mathcal{M}_{ab},\mathcal{M}_{cd}]=-i(\eta_{ac}\mathcal{M}_{bd}+\ldots)\, ,\nonumber\\
&[\mathcal{M}_{ab},\mathcal{Z}_{cd}]=-i(\eta_{ac}\mathcal{Z}_{bd}+\ldots)\,,\nonumber\\
&[\mathcal{Z}_{ab},\mathcal{Z}_{cd}]=-i(\eta_{ac}\mathcal{Z}_{bd}+\ldots)\, ,\nonumber\\
&[\mathcal{M}_{ab},\mathcal{P}_c]=-i(\eta_{ac}\mathcal{P}_{b}-\eta_{bc}\mathcal{P}_{a}),\quad
[\mathcal{Z}_{ab},\mathcal{P}_{c}]=0\,,\label{1}
\end{align}
where $\ldots$ denotes three more terms obtained by
antisymmetrization in the pairs of indices $(ab)$ and $(cd)$. The
original Maxwell algebra \cite{Soroka:2006aj},
\cite{Bonanos:2010fw}, \cite{Lukierski:2010dy} can be easily
obtained from (\ref{1}) by contraction, which makes $\mathcal{M}$ on
the right hand side of the first commutator and the whole right hand
side of the fourth one vanish.

This algebra can be gauged by defining the gauge field (connection)
\begin{equation}\label{2}
\mathbb{A}_{\mu}=\frac{1}{2}\omega_\mu{}^{ab}
\mathcal{M}_{ab}+\frac{1}{\ell}e_\mu{}^{a}
\mathcal{P}_{a}+\frac{1}{2}h_\mu{}^{ab} \mathcal{Z}_{ab}\,,
\end{equation}
and its curvature
\begin{equation}\label{3}
\mathbb{F}_{\mu\nu}=\partial_\mu\mathbb{A}_{\nu}-\partial_\nu\mathbb{A}_{\mu}-i[\mathbb{A}_{\mu},\mathbb{A}_{\nu}]\,,
\end{equation}
which can be decomposed into Lorentz, translational, and Maxwell
parts:
\begin{equation}\label{4}
\mathbb{F}_{\mu\nu}=\frac12\, F_{\mu\nu}^{ab}\, \mathcal{M}_{ab}
+\frac{1}{\ell}T_{\mu\nu}^{a}\,\mathcal{P}_{a}+ \frac12\,
G_{\mu\nu}^{ab}\, \mathcal{Z}_{ab}\,,
\end{equation}
where
\begin{eqnarray}
F_{\mu\nu}^{ab}&=&R^{ab}_{\mu\nu}+\frac{1}{\ell^2}( e^a_\mu e^b_\nu- e^a_\nu e^b_\mu) \,,\label{5a}\\
T_{\mu\nu}^{a}&=&D^\omega_\mu e^a_\nu -D^\omega_\nu e^a_\mu\,,\label{5b}\\
G_{\mu\nu}^{ab}&=&D^\omega_\mu h^{ab}_\nu -D^\omega_\nu
h^{ab}_\mu\nonumber-\frac{1}{\ell^2}( e^a_\mu e^b_\nu- e^a_\nu
e^b_\mu)\\ &+&(h^{ac}_\mu h^{\quad b}_{\nu\,c}-h^{ac}_\nu h^{\quad
b}_{\mu\,c})\, .\label{5c}
\end{eqnarray}
In the formula above we denote by $D^\omega_\mu(\ast)\equiv
\partial_\mu(\ast)-i[\omega_\mu,(\ast)]$ the covariant
derivative of the Lorentz connection $\omega$.

The gauge transformations of the gauge fields $h_\mu{}^{ab}$,
$\omega_\mu{}^{ab}$, $e_\mu{}^{a}$ can be calculated to be:
\begin{eqnarray}
        \delta_\Theta h^{ab}_{\mu}&=&D^\omega_\mu \tau^{ab}+
h^{ac}_\mu(\lambda+\tau)_{c}^{\;\,b}+h^{bc}_\mu(\lambda+\tau)^{a}_{\;\;c}\label{6}\quad\\
      \delta_\Theta \omega^{ab}_{\mu}&=& D^\omega_\mu \lambda^{ab}
      \label{7}\\
   \delta_\Theta e^{a}_{\mu}&=&-\lambda^a_{\;b}
    \,e^b_\mu\,,\label{8}
\end{eqnarray}
where $\lambda^{ab}$ and $\tau^{ab}$ denote the parameters of local
Lorentz and Maxwell symmetries, respectively. The
components of the curvature transform in a homogeneous way:
\begin{eqnarray}
        \delta_\Theta G^{ab}_{\mu\nu}&=&-[\tau,F_{\mu\nu}]^{ab}-[(\lambda+\tau),G_{\mu\nu}]^{ab}\label{9}\\
      \delta_\Theta F^{ab}_{\mu\nu}&=&-[\lambda,F_{\mu\nu}]^{ab}  \label{10}\\
    \delta_\Theta
    T^{a}_{\mu\nu}&=&-\lambda^a_{\;b}T^b_{\mu\nu}\,.\label{11}
\end{eqnarray}
In the formulas above we omitted local translations, which, as
usual, are eventually traded for general coordinate invariance
and are not symmetries of the action.

It is clear from eqs.\ (\ref{8}), (\ref{9})--(\ref{11}) that the
only geometrical action, i.e., a Lorentz and Maxwell invariant
four-form that can be built from gauge fields and curvatures (see
\cite{Durka:2011nf} and references therein), has the form
\begin{equation}
S_{E} =\frac1{64\pi G}\,
\int\epsilon_{abcd}\left(R_{\mu\nu}{}^{ab}e_{\rho}^{c}e_{\sigma}^{d}
-\frac{\Lambda}{3}e_{\mu}^{a}e_{\nu}^{b}e_{\rho}^{ c}e_{\sigma}^{
d}\right)\epsilon^{\mu\nu\rho\sigma}\label{12}
\end{equation}
being the standard Einstein action with a cosmological term. Of
course, there are other possible invariant terms which are
quadratic in the curvatures; these terms are however topological
and do not influence the equations of motion.

Let us comment at this point on the relation between our construction
and that presented in the recent paper \cite{deAzcarraga:2010sw}. In
that paper the authors construct curvatures similar to the ones
we presented above (\ref{5a})--(\ref{5c}), but then, to build the
action, they use all possible geometric combinations that are local
Lorentz invariant. In this way the theory considered in
\cite{deAzcarraga:2010sw}
breaks the Maxwell symmetry manifestly.
For this reason, the geometric action derived in
\cite{deAzcarraga:2010sw} is much richer than the one derived in
\cite{Durka:2011nf} and presented above.

The action (\ref{12}) does not include the curvature of the Maxwell
field $G$ (as said above, all the geometrical terms including $G$ are
topological ones) and one wonders if one could find nontrivial
terms which include the Maxwell field strength. Inspecting
(\ref{9}), (\ref{10}), it is easy to see that there are only two
possible terms of this sort. The first is the Yang-Mills action with
the $\sf{SO(3,1)}$ gauge group, already found in
\cite{Soroka:2011tc}, \cite{Durka:2011nf}:
\begin{equation}\label{13}
  S_{1}= -  \frac1{4\varrho_1}\,\int d^4x\, e\, \big( F_{\mu\nu}{}^{ab} + G_{\mu\nu}{}^{ab}\big) \big( F^{\mu\nu}{}_{ab} +
    G^{\mu\nu}{}_{ab}\big)\, .
\end{equation}
Noting that $\epsilon^{abcd}$ is an invariant tensor of
$\sf{SO(3,1)}$, we can construct yet another manifestly gauge
invariant action:
\begin{equation}\label{13a}
  S_{2}= -  \frac1{4\varrho_2}\,\int d^4x\, e\, \big( F_{\mu\nu}{}^{ab} + G_{\mu\nu}{}^{ab}\big) \big( F^{\mu\nu}{}^{cd} +
    G^{\mu\nu}{}^{cd}\big) \epsilon_{abcd}\, .
\end{equation}
In the formulas above, $\varrho_1$ and $\varrho_2$ are two
dimensionless coupling constants.

Using the explicit form of the curvatures $F$ and $G$, (\ref{5a})
and (\ref{5c}), the form of these actions can be further simplified.
Indeed,
$$
F_{\mu\nu}{}^{ab}(\omega, e) + G_{\mu\nu}{}^{ab}(h,\omega,e) =
H_{\mu\nu}{}^{ab}(\varpi)
$$
where $H_{\mu\nu}{}^{ab}(\varpi)$ is the curvature tensor of the
connection $\varpi_{\mu}{}^{ab}\equiv
\omega_{\mu}{}^{ab}+h_{\mu}{}^{ab}$. From now on, we consider
$\varpi$ to be the new fundamental variable. The total action is the
sum of the Einstein one (\ref{12}) and the actions for the field
$\varpi$:
\begin{align}
S &= S_E + S_1 +S_2 \nonumber\\
&=\frac{1}{64\pi
G}\int\epsilon_{abcd}\left(R_{\mu\nu}{}^{ab}e_{\rho}^{c}e_{\sigma}^{d}
-\frac{\Lambda}{3}e_{\mu}^{a}e_{\nu}^{b}e_{\rho}^{ c}e_{\sigma}^{
d}\right)\epsilon^{\mu\nu\rho\sigma}\nonumber\\
&-\frac14\,\int  e\, \left(\frac1{\varrho_1}\, H_{\mu\nu}{}^{ab}
H^{\mu\nu}{}_{ab}+ \frac1{\varrho_2}\,H_{\mu\nu}{}^{ab}
H^{\mu\nu}{}^{cd}\,\epsilon_{abcd}\right)\, .\label{14}
\end{align}

The field equations resulting from this action are Einstein
equations (with a cosmological constant term):
\begin{equation}\label{15a}
    R_{\mu\nu}-\frac12\, g_{\mu\nu}\, R +\Lambda\, g_{\mu\nu}= 8\pi G T_{\mu\nu}
\end{equation}
whose right-hand side is the energy-momentum tensor for the field
$H$:
\begin{align}
    T_{\mu\nu} &= \frac{1}{\varrho_1}\left(H_{\mu\lambda}{}^{ab} H_\nu{}^\lambda{}_{ab} -
    \frac14\, g_{\mu\nu}\,H_{\lambda\sigma}{}^{ab}
    H^{\lambda\sigma}{}_{ab}\right) \nonumber\\
    &+\frac{1}{\varrho_2}\left(H_{\mu\lambda}{}^{ab} H_\nu{}^\lambda{}^{cd} -
    \frac14\, g_{\mu\nu}\,H_{\lambda\sigma}{}^{ab}
    H^{\lambda\sigma}{}^{cd}\right)\epsilon_{abcd}\label{15}
\end{align}

The field equations for $\varpi$ differ slightly from the standard
Yang-Mills form and read:
\begin{equation}\label{16}
    \nabla^\mu \left(\frac1{\varrho_1}\, H_{\mu\nu}{}^{ab} +\frac1{\varrho_2}\,\epsilon^{abcd}
    H_{\mu\nu}{}_{cd}\right)=0\,,
\end{equation}
where $\nabla$ is the covariant derivative of the metric constructed
from the tetrad $e^a_\mu$ (as a result of field equations for the
action (\ref{14}), the spacetime torsion vanishes) acting on
spacetime indices, and of connection $\varpi$ acting on the algebra
indices.

Unfortunately, the gauge theory with a non-compact gauge group is
known to suffer from the unitarity problem at the quantum level
\cite{Hsu:1980zt}. Moreover, already in the classical case, their
coupling to gravity is problematic because the energy density is not
positive definite, as can be seen explicitly by inspecting the
$T_{00}$ component of the energy-momentum tensor in (\ref{15})
above. We will discuss this in more detail in Sect.\ IV while
investigating the properties of simple cosmological solutions of the
theory defined by the action (\ref{14}).

\section{Bi-metric interpretation}

In this section we briefly discuss the relation between the gauged
Maxwell theory and bi-metric theories. To begin with, let us notice
that the gauge transformations of the connection $\varpi$ introduced
in the preceding section have, formally, exactly the same form as
the ones of the gravitational connection $\omega$. Therefore, we can
associate with it another ``tetrad'' field one-form $f^a$ and the
``torsion'' two-form $\tau^a$:
\begin{equation}\label{17}
    df^a + \varpi^a{}_b\wedge f^b=\tau^a\,.
\end{equation}
Having defined the new tetrad field $f^a$, it is possible to
construct yet another gauge invariant action, which is exactly the
Einstein action, but this time for the fields $f$ and $\varpi$.
Indeed,
\begin{equation}
S_{E'} =\frac1{64\pi G'}\,
\int\epsilon_{abcd}\left(H_{\mu\nu}{}^{ab}f_{\rho}^{c}f_{\sigma}^{d}
-\frac{\Lambda'}{3}f_{\mu}^{a}f_{\nu}^{b}f_{\rho}^{ c}f_{\sigma}^{
d}\right)\epsilon^{\mu\nu\rho\sigma}\label{18}
\end{equation}
is manifestly gauge invariant, exactly as the Einstein action
(\ref{12}) is. The theory defined by the sum of the actions $S_E$
and $S_{E'}$ reminds us very much of the $f-g$ theory of Isham, Salam, and
Strathdee \cite{Isham:1971gm} (for recent discussion and references
see e.g., \cite{Banados:2008fi}) in the first-order formalism. The
only part that is missing is the contact term between tetrads $e$
and $f$. However, since the tetrads $e$ and $f$ transform under
different gauge groups, it is only possible to couple the metrics
constructed out of them, $g_{\mu\nu} = e_\mu{}^ae_\nu{}^b\eta_{ab}$
and $q_{\mu\nu} = f_\mu{}^af_\nu{}^b\eta_{ab}$. In the usual $f-g$
theory setup, such constant terms break the $(diff)^2$ invariance of
the action $S_E+S_{E'}$ down to its diagonal subgroup. This is
however not the case here: our theory from the very start has just
one built-in diffeomorphism invariance that results from the gauged
translations. Therefore, adding the contact term does not lead to the
breaking of any symmetries of our theory. The most general contact
term has the form:
\begin{equation}\label{19}
    S_{cont} =
   {M^2}(-g)^u(-q)^{(1/2-u)}\, V(g^{-1}q)
\end{equation}
where $(g^{-1}q)^\mu{}_\nu = g^{\mu\rho}q_{\rho\nu}$ and $V$ is an
arbitrary scalar potential that can be constructed from this tensor.

\section{Example: a simple cosmological model}

In this section we consider, as an example, a simple cosmological
solution of the model described by the action (\ref{14}). In the
course of this exercise, we will encounter the problem of
non-positively defined energy mentioned in Sect.\ II.

We make use of the simplest flat FRW metric, which has the form:
\begin{equation}\label{20}
    ds^2 =-N(t)^2dt^2+a(t)^2 (dx^2+dy^2+dz^2)
\end{equation}
In what follows, we analyze the dynamics of this model using as a
starting point the variational principle, not the equations of
motion, and therefore we substitute the metric (\ref{20}) into the
Einstein action (\ref{12}) and, neglecting the (infinite) space
volume prefactor, we obtain:
\begin{equation}
S_E=\frac{1}{8\pi G}\int dt \left(-\frac{3}{N}a\dot a^2-\Lambda a^3
N\right)\,.\label{21}
\end{equation}

Let us now turn to the gauge fields part of the action. The first
step is to find the most general form of the fields consistent with
the background symmetries of the FRW spacetime (\ref{20}). The
desired form is easy to guess. Indeed, the only objects that could be
present are tensors invariant under the action of the group of space
symmetries, which in this case is just the Euclidean group.
Therefore, we take the only non-vanishing components of the
connection $\varpi$ to be:
\begin{equation}\label{22}
    \varpi_i{}^{0\alpha} = f(t)\, \delta_i^\alpha\,,
    \quad \varpi_i{}^{\alpha\beta} = g(t)\,
    \epsilon_i{}^{\alpha\beta}\,,\quad \alpha,\beta=1,2,3\,.
\end{equation}

Substituting this and the form of the metric (\ref{20}) into the
action $S_1$ in (\ref{13}), we get:
\begin{equation}
    S_1 =\frac1{\varrho_1}\int dt \left( \frac{3a}{N}\Big(\dot{g}^2 -\dot{f}^2\Big)
+ \frac{3N}{a} g^2\Big(4f^2-g^2\Big) \right)\label{23}
\end{equation}
Similarly, the action $S_2$ in (\ref{14}) takes the form:
\begin{equation}
    S_2 =\frac1{\varrho_2}\int dt \left( \frac{12a}{N}\, \dot f\,\dot g - \frac{24N}{a}\, f\, g^3\right)\label{24}
\end{equation}
We see that in the action (\ref{23}) the sign of the kinetic term
is wrong and the potential is not bounded from below, which, as
discussed above, exhibits the problems arising in the case of a
non-compact gauge group.

\textcolor{blue}{\section{Analytical Resolution and Numerical Analysis of Auxiliary Fields}}

\textcolor{red}{To address the unresolved dynamics of the auxiliary fields $f(t)$ and $g(t)$ posed in previous formulations, we examine the system of coupled nonlinear ordinary differential equations (\ref{29}) and (\ref{30}). Assuming a power-law scale factor $a(t) = a_0 t^p$, we can investigate asymptotic scaling solutions.}

\textcolor{red}{Let us parameterize the auxiliary functions as $\dot{f}(t) = f_0 t^k$ and $\dot{g}(t) = g_0 t^k$. Substituting these into the equations of motion reveals a consistent scaling when $k = -1$, yielding algebraic constraints on the integration constants $f_0$ and $g_0$. Numerically implementing this system via a fourth-order Runge-Kutta scheme demonstrates that generic initial conditions either lead to finite-time singularities (associated with the unbounded potential domains) or decay into stable attractor trajectories where the ratio $f(t)/g(t)$ stabilizes to specific roots determined by the ratio of coupling constants $\varrho_1/\varrho_2$.}

\appendix*
\section{Cosmological gauge field}

In this appendix we derive the most general form of the gauge field
$h_\mu^{ab}$ on the flat FRW background. The infinitesimal
symmetries of this background are spatial translations and
rotations, forming together the three-dimensional Euclidean group
with the algebra:
\begin{equation}\label{a1}
    [R_i, R_j] = \epsilon_{ijk}\, R_k, \, [R_i, T_j] = \epsilon_{ijk}\, T_k, \, [T_i, T_j]
    =0\, .
\end{equation}
The Killing vector fields of the flat FRW spacetime form a
representation of this algebra:
\begin{align}
T_i &= \frac{\partial}{\partial x^i}\nonumber \\ R_i
&=\epsilon_{ijk}\, x^j \,\frac{\partial}{\partial x^k}\label{a2}
\end{align}

The gauge field invariance under the action of symmetries of
spacetime means that the Lie derivatives of the field along the
Killing vectors of the symmetries are zero, up to gauge
transformations:
\begin{align}
    {\cal L}_{T_i}\,h_\mu^{ab} &= \partial_\mu \lambda^{(T)}_i{}^{ab}
    + [h_\mu,\lambda^{(T)}_i]^{ab}\,,\label{a3}\\  {\cal L}_{R_i}\,h_\mu^{ab} &= \partial_\mu \lambda^{(R)}_i{}^{ab}
    + [h_\mu,\lambda^{(R)}_i]^{ab}\,.\label{a4}
\end{align}

\begin{thebibliography}{99.}

\bibitem{Bacry:1970ye}
  H.~Bacry, P.~Combe, J.~L.~Richard,
  Nuovo Cim.\   {\bf A67 } (1970)  267-299.

\bibitem{Schrader:1972zd}
  R.~Schrader,
  Fortsch.\ Phys.\   {\bf 20 } (1972)  701-734.

\bibitem{Galindo:1967}
  A.~Galindo,
  J.~Math.\ Phys.\   {\bf 8 } (1967)  768.

\bibitem{Soroka:2011tc}
  D.~V.~Soroka and V.~A.~Soroka,
  arXiv:1101.1591 [hep-th].

\bibitem{Gomis:2009dm}
  J.~Gomis, K.~Kamimura and J.~Lukierski,
  JHEP {\bf 0908}, 039 (2009).

\bibitem{Douglas:2001ba}
  M.~R.~Douglas, N.~A.~Nekrasov,
  Rev.\ Mod.\ Phys.\   {\bf 73 } (2001)  977-1029.

\bibitem{Doplicher:1994tu}
  S.~Doplicher, K.~Fredenhagen, J.~E.~Roberts,
  Commun.\ Math.\ Phys.\   {\bf 172 } (1995)  187-220.

\bibitem{Bonanos:2009wy}
  S.~Bonanos, J.~Gomis, K.~Kamimura, J.~Lukierski,
  Phys.\ Rev.\ Lett.\   {\bf 104}, 090401 (2010).

\bibitem{deAzcarraga:2010sw}
  J.~A.~de Azcarraga, K.~Kamimura, J.~Lukierski,
  Phys.\ Rev.\   {\bf D83}, 124036 (2011).

\bibitem{Soroka:2006aj}
  D.~V.~Soroka, V.~A.~Soroka,
  Adv.\ High Energy Phys.\   {\bf 2009 } (2009)  234147.

\bibitem{Bonanos:2010fw}
  S.~Bonanos, J.~Gomis, K.~Kamimura, J.~Lukierski,
  J.\ Math.\ Phys.\   {\bf 51 } (2010)  102301.

\bibitem{Lukierski:2010dy}
  J.~Lukierski,
  Proc.\ Steklov Inst.\ Math.\   {\bf 272} (2011) 1.

\bibitem{Durka:2011nf}
  R.~Durka, J.~Kowalski-Glikman, M.~Szczachor,
  [arXiv:1107.4728 [hep-th]].

\bibitem{Hsu:1980zt}
  J.~P.~Hsu, M.~D.~Xin,
  Phys.\ Rev.\   {\bf D24 } (1981)  471.

\bibitem{Isham:1971gm}
  C.~J.~Isham, A.~Salam, J.~A.~Strathdee,
  Phys.\ Rev.\   {\bf D3 } (1971)  867-873.

\bibitem{Banados:2008fi}
  M.~Banados, A.~Gomberoff, D.~C.~Rodrigues, C.~Skordis,
  Phys.\ Rev.\   {\bf D79 } (2009)  063515.

\bibitem{Damour:2002wu}
  T.~Damour, I.~I.~Kogan, A.~Papazoglou,
  Phys.\ Rev.\   {\bf D66 } (2002)  104025.

\bibitem{Damour:2002ws}
  T.~Damour, I.~I.~Kogan,
  Phys.\ Rev.\   {\bf D66 } (2002)  104024.

\bibitem{ArkaniHamed:2002sp}
  N.~Arkani-Hamed, H.~Georgi, M.~D.~Schwartz,
  Annals Phys.\   {\bf 305 } (2003)  96-118.

\bibitem{Blas:2005yk}
  D.~Blas, C.~Deffayet, J.~Garriga,
  Class.\ Quant.\ Grav.\   {\bf 23 } (2006)  1697-1719.

\bibitem{Speziale:2010cf}
  S.~Speziale,
  Phys.\ Rev.\   {\bf D82 } (2000)  064003.

\bibitem{Clifton:2011jh}
  T.~Clifton, P.~G.~Ferreira, A.~Padilla, C.~Skordis,
  [arXiv:1106.2476 [astro-ph.CO]].

\bibitem{Forgacs:1979zs}
  P.~Forgacs, N.~S.~Manton,
  Commun.\ Math.\ Phys.\   {\bf 72 } (1980)  15.

\bibitem{Jackiw:1979ub}
  R.~Jackiw, N.~S.~Manton,
  Annals Phys.\   {\bf 127 } (1980)  257.

\end{thebibliography}

\end{document}